\documentclass[a4wide,11pt]{article}
\usepackage[utf8]{inputenc}
\usepackage{CJKutf8}
\usepackage{graphicx}
\usepackage{url}
\usepackage{float}
\usepackage{csvsimple}
\usepackage[pdfborder=0 0 0]{hyperref}
\usepackage{longtable}
\usepackage{fancyhdr}
\pagestyle{fancy}
\usepackage{graphicx}

\begin{document}
\begin{CJK}{UTF8}{bsmi}

\title{Standard Grimoire Report \\
  PROJECT-NAME Project \\
      2016-QUARTER \\
        ~~\\~~\\  \includegraphics[scale=.35]{logo}       \\}




\maketitle


\newpage
~~\\~\\~\\~\\
%\begin{tabular}{p{5cm}p{10cm}}
%~  & Daniel Izquierdo Cort\'azar \\
%       & Chief Data Officer  \\
%       & dizquierdo@bitergia.com \\

%\end{tabular}

~~\\~~\\

~~\\
~~\\

\begin{tabular}{p{5cm}p{10cm}}

\end{tabular}

~~\\
~~\\

This report would not exist without the effort of the people involved in the development of the \href{http://grimoirelab.github.io}{Grimoire toolset}.

~~\\
~~\\
~~\\
~~\\
~~\\
~~\\

\begin{flushright}
  (cc) 2016 Bitergia. Some rights reserved. \\
  This work licensed under Creative Commons
  Attribution-ShareAlike 4.0 Unported License. \\
  To view a copy of full license, see
  \url{http://creativecommons.org/licenses/by-sa/4.0}, \\
  or write to
  Creative Commons, 559 Nathan Abbott Way, Stanford,
  California 94305, USA. \\ \ 
\end{flushright}


%\rhead{Bitergia Report}
%\rhead{\includegraphics[scale=.15]{openstack.eps}\includegraphics[scale=.15]{logo.eps}}
\rhead{\includegraphics[scale=.15]{logo}}
%\rhead{\includegraphics[scale=.15]{logo.eps}}


\newpage

\begin{center}
\textbf{Executive Summary}
\end{center}

This report provides a quantitative analysis of the current and past situation of the PROJECT-NAME project. All the data presented in it is based on information retrieved from the software development repositories of the project. The analysis includes a summary of the general situation of the project, and specific analysis of some of its development processes and communication channels. Data from previous periods is also shown for comparison.

\newpage

\setcounter{tocdepth}{4}
\setcounter{secnumdepth}{4}
\tableofcontents

\newpage

\section{Project overview} \label{overview}

The report looks at activities across the PROJECT-NAME community during
2016-QUARTER (DATEFRAME), comparing it to previous period of analysis.

\begin{table}[H]
    \centering
    \begin{tabular}{c|c|c|}
		% specify table head
	    \bfseries Data source & \bfseries Activity last quarter & \bfseries Change (wrt to prev. quarter)
	    % use head of csv as column names
	    \csvreader[head to column names]{data/data_source_evolution.csv}{}
	    {\\\datasource & \netvalues ~ \metricsnames & \relativevalues\% }
    \end{tabular}
    \caption{Activity during the last period of analysis and its evolution}
    \label{tab:_summary}
\end{table}

Table~\ref{tab:_summary} shows development activity for each of the analyzed
data sources. The activity column displays information about the net
activity numbers, while the Change column displays information
about the relative difference with respect to the previous period of
analysis.\\

%TODO: add list of repositories analyzed per quarter

%\section{Development and Conversations}

%The mean number of developers per month

%\begin{figure}[H]
%    \centering
%    \includegraphics[scale=.35]{figs/authors_month.eps}
%\caption{Number of developers per month in mean for each quarter}
%\end{figure}

In order to provide some insight into the software development process of PROJECT-NAME community, several metrics on GitHub Issues and Pull Requests are presented below. In terms of GitHub terminology, these metrics are calculated as follows:
\begin{itemize}
	\item Review Efficiency Index (REI), measured as the number of closed pull requests out of the new ones in a given period.
	\item Time to Merge (TTM), measured as the time since a pull request is submitted until this is closed.
	\item Backlog Management Index (BMI), measured as the number of closed tickets out of open tickets in a given period.
	\item Time to Close issues (TTC), the time since an issue was opened till this is closed.
\end{itemize}


\begin{table}[H]
    \centering
    \begin{tabular}{c|c|c|l}%
    \bfseries  REI  & \bfseries BMI & \bfseries TTM  & \bfseries TTC
    \csvreader[head to column names]{data/efficiency.csv}{}% use head of csv as column names
    {\\\bmipr & \bmitickets & \daystocloseprmedian ~ days & \daystocloseticketmedian ~ days}
    \end{tabular}
    \caption{Closed out of opened pull requests(REI), closed out of opened issues (BMI), median time to merge pull requests (TTM) and median time to close issues(TTC)}
\end{table}

Gerrit data offers an interesting view on how software development process of PROJECT-NAME community is working from the point of view of efficiency and process. To analyze this, some metrics based on Gerrit changesets and reviews are presented below. In terms of Gerrit terminology, these metrics are calculated as follows:
\begin{itemize}
	\item Review Efficiency Index (REI), measured as the number of closed changesets out of the new ones in a given period.
	\item Time to Merge (TTM), measured as the time since a review is submitted until this is closed.
\end{itemize} 


\begin{table}[H]
    \centering
    \begin{tabular}{c|c|}%
    \bfseries  REI  & \bfseries TTM % specify table head
    \csvreader[head to column names]{data/efficiency.csv}{}% use head of csv as column names
    {\\\bmireviews  & \daystomergereviewmedian ~ days }
    \end{tabular}
    \caption{Closed out of opened changesets (REI) and median time to merge (TTM)}
\end{table}

The rest of the document is divided into three sections with information from the last periods: 

\begin{itemize}
\item Activity: focused on contributions. 
\item Community: focused on contributors.
\item Process: focused on efficiency and timing.
\end{itemize}

\section{Activity} \label{activity}
This section covers contributions in the different data sources.\\

The bar chart below shows the evolution of the number of commits and authors in Git through time, grouped by quarters.

\begin{tabular}{p{9cm} p{5cm}}
	\vspace{0pt} 
	\hspace*{-6cm}  
	\includegraphics[scale=.75]{activity/git_commits_git_authors.eps}
	& 
	\vspace{0pt}
	\begin{tabular}{l|r|r|}
		% specify table head
		\bfseries Period & \bfseries Commits & \bfseries Authors
		% use head of csv as column names
		\csvreader[head to column names]{activity/git_commits_git_authors.csv}{}
		{\\\Date & \commits & \authors}
	\end{tabular}
\end{tabular}

%The total number of contributors divided into three sets (core,
%regular and casual\footnote{Contributing developers are characterized
%as core, regular and casual depending on their activity in the git
%repositories. The classification is built by sorting contributors by
%their total number of commits; we sum the total commits per each
%individual contributors: the individuals whose commits sum up to 80\%
%of the total number of commits in the quarter are the core
%contributors in that quarter. The regular contributors are those whose
%commits sum up to 95\% of the total. The others are the casual
%contributors.}) follow a similar pattern.

%\begin{figure}[H]
%    \centering
%    \includegraphics[scale=.35]{figs/onion.eps}
%    \caption{Evolution during the last quarters of core, regular and casual developers (based on git activity)}
%\end{figure}

%\begin{table}[H]
%    \centering
%    \begin{tabular}{l|r|r|r|}%
%    \bfseries Period & \bfseries Core & \bfseries Regular & \bfseries Occasional% specify table head
%    \csvreader[head to column names]{data/onion_model.csv}{}% use head of csv as column names
%    {\\\labels & \core & \regular & \occasional}
%    \end{tabular}
%    \caption{Characterization of developers by their total contribution to the project}
%\end{table}

Following chart shows activity in GitHub data source. It is based on comparing the number of opened and closed reviews grouped by quarters. A review is equivalent to a pull request in GitHub terminology.

\begin{tabular}{p{7cm} p{5cm}}
	\vspace{0pt} 
	\includegraphics[scale=.35]{figs/github_submitted_github_closed_general.eps}
	& 
	\vspace{0pt}
	\begin{tabular}{l|r|r|}%
		\bfseries Period & \bfseries Opened & \bfseries Closed % specify table head
		\csvreader[head to column names]{data/github_submitted_github_closed_general.csv}{}% use head of csv as column names
		{\\\labels & \submitted & \closed}
	\end{tabular}
\end{tabular}
The following chart shows activity in Gerrit data source. It is based on comparing the number of opened and closed reviews grouped by quarters.

\begin{tabular}{p{8cm} p{5cm}}
	\vspace{0pt} 
	\hspace*{-5cm}  
	\includegraphics[scale=.75]{activity/gerrit_submitted_gerrit_closed.eps}
	& 
	\vspace{0pt}
	\begin{tabular}{l|r|r|}%
		\bfseries Period & \bfseries Opened & \bfseries Closed % specify table head
		\csvreader[head to column names]{activity/gerrit_submitted_gerrit_closed.csv}{}% use head of csv as column names
		{\\\Date & \submitted & \closed}
	\end{tabular}
\end{tabular}
The following chart shows the evolution of the community in terms of people sending emails, grouped by quarters for comparison. 

\begin{tabular}{p{7cm} p{5cm}}
    \vspace{0pt} 
    \includegraphics[scale=.35]{figs/mls_emails_senders.eps}
    & 
    \vspace{0pt}
    \begin{tabular}{l|l}%
    \bfseries Period & \bfseries People % specify table head
    \csvreader[head to column names]{data/mls_emails_senders.csv}{}% use head of csv as column names
    {\\\labels & \emailssenders}
    \end{tabular}
\end{tabular}



\section{Community} \label{community}
This section tries to help us to understand the evolution of PROJECT-NAME community by looking at active contributors and organizations in the last period of analysis, compared to previous ones.\\

The bar chart below shows the evolution of the number of commits and authors in Git through time, grouped by quarters.

\begin{tabular}{p{9cm} p{5cm}}
	\vspace{0pt} 
	\hspace*{-6cm}  
	\includegraphics[scale=.75]{activity/git_commits_git_authors.eps}
	& 
	\vspace{0pt}
	\begin{tabular}{l|r|r|}
		% specify table head
		\bfseries Period & \bfseries Commits & \bfseries Authors
		% use head of csv as column names
		\csvreader[head to column names]{activity/git_commits_git_authors.csv}{}
		{\\\Date & \commits & \authors}
	\end{tabular}
\end{tabular}

%The total number of contributors divided into three sets (core,
%regular and casual\footnote{Contributing developers are characterized
%as core, regular and casual depending on their activity in the git
%repositories. The classification is built by sorting contributors by
%their total number of commits; we sum the total commits per each
%individual contributors: the individuals whose commits sum up to 80\%
%of the total number of commits in the quarter are the core
%contributors in that quarter. The regular contributors are those whose
%commits sum up to 95\% of the total. The others are the casual
%contributors.}) follow a similar pattern.

%\begin{figure}[H]
%    \centering
%    \includegraphics[scale=.35]{figs/onion.eps}
%    \caption{Evolution during the last quarters of core, regular and casual developers (based on git activity)}
%\end{figure}

%\begin{table}[H]
%    \centering
%    \begin{tabular}{l|r|r|r|}%
%    \bfseries Period & \bfseries Core & \bfseries Regular & \bfseries Occasional% specify table head
%    \csvreader[head to column names]{data/onion_model.csv}{}% use head of csv as column names
%    {\\\labels & \core & \regular & \occasional}
%    \end{tabular}
%    \caption{Characterization of developers by their total contribution to the project}
%\end{table}

The following chart shows the evolution of the community in terms of people sending emails, grouped by quarters for comparison. 

\begin{tabular}{p{7cm} p{5cm}}
    \vspace{0pt} 
    \includegraphics[scale=.35]{figs/mls_emails_senders.eps}
    & 
    \vspace{0pt}
    \begin{tabular}{l|l}%
    \bfseries Period & \bfseries People % specify table head
    \csvreader[head to column names]{data/mls_emails_senders.csv}{}% use head of csv as column names
    {\\\labels & \emailssenders}
    \end{tabular}
\end{tabular}



\section{Process} \label{process}
This section intends to show the evolution of efficiency and timing when dealing with tasks related with code review processes.\\
	
To measure the efficiency of the process of software development in GitHub first we focus on BMI, that was defined in section \ref{overview}. Chart below shows evolution of BMI by quarters, allowing us to check how pull requests are being managed by the community in the last quarter compared to previous ones.

\begin{tabular}{p{7cm} p{5cm}}
	\vspace{0pt} 
	\includegraphics[scale=.35]{figs/github_prs_bmipr_general.eps}
	& 
	\vspace{0pt}
	\begin{tabular}{l|l}%
		\bfseries Period & \bfseries Closed/Subm. % specify table head
		\csvreader[head to column names]{data/github_prs_bmipr_general.csv}{}% use head of csv as column names
		{\\\labels & \bmipr}
	\end{tabular}
\end{tabular}

Besides, next chart deals with timing. It shows mean and median times to merge pull requests in GitHub (TTM, defined in section \ref{overview})--in days--for last quarter compared to previous ones.


\begin{tabular}{p{7cm} p{5cm}}
	\vspace{0pt} 
	\includegraphics[scale=.35]{figs/github_prs_days_to_close_pr_avg_github_prs_days_to_close_pr_median_general.eps}
	& 
	\vspace{0pt}
	\begin{tabular}{l|r|r|}%
		\bfseries Period & \bfseries Median & \bfseries Mean % specify table head
		\csvreader[head to column names]{data/github_prs_days_to_close_pr_avg_github_prs_days_to_close_pr_median_general.csv}{}% use head of csv as column names
		{\\\labels & \daystocloseprmedian & \daystoclosepravg}
	\end{tabular}
\end{tabular}

The following chart shows activity in Gerrit data source. It is based on comparing the number of opened and closed reviews grouped by quarters.

\begin{tabular}{p{8cm} p{5cm}}
	\vspace{0pt} 
	\hspace*{-5cm}  
	\includegraphics[scale=.75]{activity/gerrit_submitted_gerrit_closed.eps}
	& 
	\vspace{0pt}
	\begin{tabular}{l|r|r|}%
		\bfseries Period & \bfseries Opened & \bfseries Closed % specify table head
		\csvreader[head to column names]{activity/gerrit_submitted_gerrit_closed.csv}{}% use head of csv as column names
		{\\\Date & \submitted & \closed}
	\end{tabular}
\end{tabular}


\appendix{}

\section{Metrics Definitions} \label{metrics}
\label{chap:metrics_definitions}

\begin{itemize}
\item Commit: this is defined as the action(s) that performes a change in the source code. Bots, merges and other type of automatic activity is removed from the records. In addition, when aggregating several git repositories, this metric only counts unique revisions (unique hashes found in the git repositories). In addition, all branches are aggregated to the analysis.

\item Submitted changesets: a code review is the process of peer reviewing source code changes. A submitted code is not merged to the master code of a given project till this is approved. A submitted code review is defined as any changeset submitted to the Gerrit
system. 
 
\item Authors: a developer is defined as author if she is the owner of the patchset sent for reviewing and this is merged into the source code. As previously indicated, automatic commits such bot's are removed from this analysis.

\item Efficiency closing changesets: this metric is a derivation of the Backlog Management Index as it is named as Review efficiency index (REI). As similarly used in the BMI index, this metrics measures the number of closed changesets
 out of the total number of new changesets in a given period.

\item Time to Merge: this time consists of the time between the first upload of the first changeset till the last iteration of the code review process is merged into the code. This metric is provided in number of days.

\item Developer per period: average of developers per period ignoring bots and merges.

\item Emails sent: number of emails sent by people to the several mailing lists. Bots are not registered.

\item People sending emails: number of people sending those emails ignoring bots.

\end{itemize}


%\chapter{Source code and data sources}
%\label{chap:data_sources}

%The source code of the scripts and templates used to produce this
%report are available from the GrimoireReports repository\footnote{\url{
%https://github.com/grimoirelab}}.
%Commit identified with the hash \emph{d33b6f5951662fb5bd895eee89c0b0344b09c956}. 

%select p.id, pr.data_source, pr.repository_name 
%from projects p, project_repositories pr 
%where p.project_id=pr.project_id order by p.id, pr.data_source;

%\input{repositories.tex}

\end{CJK}
\end{document}
